\section{Discussion and Limitations}
\label{sec:discussion}
\subsection{Planning and acceleration compose}
The central observation of this study is that a planner and a GPU attack different factors of the same cost. Direct dense simulation of a batch performs roughly $\sum_p G_p 2^{n_p}$ amplitude-gate updates. A GPU lowers the time of each update, while the planner changes how many updates are needed: it replaces $2^{n}$ by $2^{w}$ for every distinct cone and the number of prepared states by the number of distinct cone circuits. When cones are narrow, the second factor wins by an exponential margin and the GPU has nothing left to do: at 28 qubits the planner, on one CPU thread, spends 3.8\,ms on the laptop and 4.6 to 6.8\,ms on the cloud host, where Aer spends 4.35\,s on the laptop GPU and 0.39\,s on an H100. When cones saturate, there is nothing to remove, and the right behaviour is to hand the work to the fastest device, which is what width-based placement does: on our saturated control, it matched direct execution on the laptop GPU within 0.6\% from 24 qubits on, and on the H100 at 28. The two techniques are therefore complements rather than competitors, and the question a planner should answer is not whether it can replace the GPU but which part of the requested work needs it.

The 48-qubit intermediate-cone cases make this separation concrete. At depth eight the planner's tasks are 10--18 qubits wide: too wide for efficient serial CPU execution, but small enough for the GPU even though a full 48-qubit state is impossible. GPU placement reduces total time from 1.476\,s to 171\,ms without changing the selected plan. This is the value of treating plan choice and device placement as separate decisions.

\subsection{What a datacenter GPU changes}
\begin{table}[!t]
\caption{Fastest Direct Method's Time over the Planner's, Each Measured Within One Host: the Laptop RTX 4060 (Lap.) and a Cloud H100 PCIe (Rem.), by Number of Qubits.}
\label{tab:remote}
\centering
\footnotesize
\resizebox{\columnwidth}{!}{\begin{tabular}{llrrrrrrrr}
\toprule
 & & \multicolumn{2}{c}{16} & \multicolumn{2}{c}{20} & \multicolumn{2}{c}{24} & \multicolumn{2}{c}{28}\\
\cmidrule(lr){3-4}\cmidrule(lr){5-6}\cmidrule(lr){7-8}\cmidrule(lr){9-10}
Family & Outputs & Lap. & Rem. & Lap. & Rem. & Lap. & Rem. & Lap. & Rem.\\
\midrule
Hardware-eff. & $Z$, all & 2.05 & 1.86 & 6.34 & 1.65 & 73.3 & 8.09 & 1{,}137 & 57.0\\
Hardware-eff. & $X,Y,Z$, one & 14.0 & 11.8 & 79.5 & 31.1 & 960 & 98.3 & 15{,}279 & 1{,}203\\
QCNN-like & $Z$, all & 2.46 & 1.64 & 10.2 & 3.98 & 122 & 17.0 & 1{,}681 & 119\\
QCNN-like & $X,Y,Z$, one & 9.55 & 6.88 & 39.6 & 17.5 & 529 & 57.7 & 6{,}111 & 592\\
Re-uploading & $Z$, all & 0.36 & 0.22 & 0.74 & 0.25 & 1.00 & 0.64 & 1.00 & 0.95\\
Re-uploading & $X,Y,Z$, one & 0.50 & 0.27 & 0.83 & 0.24 & 1.01 & 0.69 & 1.00 & 0.97\\
\bottomrule
\end{tabular}
}
\end{table}
Our main GPU is a laptop device with 8\,GB of memory and about 256\,GB/s of memory bandwidth. We therefore repeated every GPU-dependent study on a cloud host with an NVIDIA H100 PCIe of 80\,GB, measuring every planner/direct ratio within that host (Table~\ref{tab:remote}). The H100 ran the 28-qubit dense circuits 11 to 18 times faster than the laptop GPU, while the host's CPU thread ran the planner about half as fast as the laptop's. The result is the shift we expected rather than a change of regime. For the local families the margins at 28 qubits shrink by a factor of 10 to 20, to 57 to 119 times with $Z$ on every wire and 592 to 1,203 times with three outputs on one wire, but the planner still wins from 16 qubits (1.6 to 1.9 times) and its advantage still grows exponentially. Gradient crossovers move by a few qubits. The full hardware-efficient Jacobian still favours the planner from 20 qubits, but by 1.3 rather than 5.7 times there and by 12 rather than 135 times at 24 qubits, where the 80\,GB device also avoids the out-of-memory failure of the adjoint method. The ten-step training loop ties with GPU adjoint at 24 qubits instead of being 13.6 times faster. On the saturated family the planner's fixed costs weigh more against a faster GPU, and placement reaches parity only at 28 qubits (0.95 to 0.97 times, against 0.22 to 0.69 between 16 and 24). An H100 also moves the memory wall from 28 to about 32 qubits in double precision, which we did not explore.

\subsection{Where the incremental contribution lies}
The stronger baselines change the interpretation of the results. Qiskit's LightCone pass already performs commutation-aware reduction, and a global MPS already shares one tensor representation across outputs. Consequently, neither cone reduction nor multi-output reuse can be the paper's novelty claim. The measured increment is the planning layer: structure-keyed analysis across a shift batch, exact grouping of the resulting bound preparations, an explicit decision about when analysis pays, and placement of the remaining tasks. The LightCone ablation isolates this layer by charging every Qiskit pass and deduplication step; the MPS experiments show that the resulting plan still has to compete with a fundamentally different representation.

The same experiments also bound that contribution. Qiskit's commutation analysis can remove gates our support traversal keeps. Exact MPS evaluates the saturated 48-qubit depth-four batch in 75.9\,ms, while LazyQC refuses it. At 128 qubits, the local-family advantage over exact MPS is 7.4--17.2 times, not the three-orders-of-magnitude margin against direct dense execution. A defensible claim is therefore not ``cones are faster than tensor networks,'' but ``a cost-aware batch planner can choose and reuse reduced state-vector work effectively in measured regimes, and its decisions can be evaluated independently of the backend.''

\subsection{Limitations}
Our conclusions hold for the workloads we measured, and several of their features favour or disfavour cones in ways a reader should keep in mind:
\begin{itemize}
\item \textbf{Workloads.} Every direct GPU result comes from the main comparison: generated families of depth four on two held-out seeds. The generalization studies widen this to two unseen families at depths 2 to 16 with five seeds and to 63 trained classifiers, but only on one CPU thread, at up to 20 qubits, and the trained models have at most 8 qubits. The LightCone comparison covers the hardware-efficient family with $Z$ outputs, and the shared-MPS study covers the generated families at depths four, six and eight. Trained models wide enough for cones to matter remain untested.
\item \textbf{Conservative cones.} The support-based traversal keeps gates that overlap but commute with the propagated observable. Qiskit's commutation-aware pass can produce narrower cones; implementing a commutation-aware analysis while retaining structure-keyed batch reuse is future work.
\item \textbf{Comparison scope.} QuWARP is a related shared-prefix planner, not a head-to-head baseline: our implementation and artifacts do not reproduce its Qrack execution surface or typed continuations. Pauli-propagation methods are discussed but not benchmarked. Tensor-network conclusions are limited to the tested quimb contraction and exact MPS configurations, not every contraction strategy or approximate tensor method.
\item \textbf{Two hosts.} Most measurements come from one laptop, and the direct GPU studies were repeated on one cloud host with an H100. Absolute times and the placement threshold depend on the host, and the analysis-budget coefficients and $w^\star$ were fitted on the laptop. The new LightCone and tensor-network studies ran on the laptop.
\item \textbf{Fixed per-task costs.} The work score has no per-task constant, so below a millisecond it can prefer a cone plan that the dense plan beats (RQ3). \textsc{LazyQC} also lowers native circuits at every collection, which costs about 25\% against a prebuilt direct circuit at 8 qubits, and its planning cache keys on parameter bindings, so it cannot amortize analysis across inference batches with new data. The policy ablation diagnoses this issue but does not tune a replacement on the held-out cases.
\item \textbf{Placement after selection.} The planner chooses a plan with a CPU-oriented score and only then places wide tasks on the GPU. Selecting plans, representations and devices jointly, or batching many narrow cone tasks into one GPU kernel, could change decisions near the crossover.
\item \textbf{Exact expectations.} The core comparisons evaluate exact expectation values. Shot noise, noisy simulation and the statistics of measurement grouping are outside the evaluated scope, and the hardware probes of Section~\ref{sec:rq4} are exploratory, with two seeds on one device. The optional sensitivity approximation is not part of the exact planner and its bound assumes an explicit numerical allowance rather than outward-rounded floating-point certification.
\end{itemize}
